\IfFileExists{stacks-project.cls}{%
\documentclass{stacks-project}
}{%
\documentclass{amsart}
}

% For dealing with references we use the comment environment
\usepackage{verbatim}
\newenvironment{reference}{\comment}{\endcomment}
%\newenvironment{reference}{}{}
\newenvironment{slogan}{\comment}{\endcomment}
\newenvironment{history}{\comment}{\endcomment}

% For commutative diagrams we use Xy-pic
\usepackage[all]{xy}

% We use 2cell for 2-commutative diagrams.
\xyoption{2cell}
\UseAllTwocells

% We use multicol for the list of chapters between chapters
\usepackage{multicol}

% This is generally recommended for better output
\usepackage{lmodern}
\usepackage[T1]{fontenc}

% For cross-file-references
\usepackage{xr-hyper}

% Package for hypertext links:
\usepackage{hyperref}

% For any local file, say "hello.tex" you want to link to please
% use \externaldocument[hello-]{hello}
\externaldocument[introduction-]{introduction}
\externaldocument[conventions-]{conventions}
\externaldocument[sets-]{sets}
\externaldocument[categories-]{categories}
\externaldocument[topology-]{topology}
\externaldocument[sheaves-]{sheaves}
\externaldocument[sites-]{sites}
\externaldocument[stacks-]{stacks}
\externaldocument[fields-]{fields}
\externaldocument[algebra-]{algebra}
\externaldocument[brauer-]{brauer}
\externaldocument[homology-]{homology}
\externaldocument[derived-]{derived}
\externaldocument[simplicial-]{simplicial}
\externaldocument[more-algebra-]{more-algebra}
\externaldocument[smoothing-]{smoothing}
\externaldocument[modules-]{modules}
\externaldocument[sites-modules-]{sites-modules}
\externaldocument[injectives-]{injectives}
\externaldocument[cohomology-]{cohomology}
\externaldocument[sites-cohomology-]{sites-cohomology}
\externaldocument[dga-]{dga}
\externaldocument[dpa-]{dpa}
\externaldocument[sdga-]{sdga}
\externaldocument[hypercovering-]{hypercovering}
\externaldocument[schemes-]{schemes}
\externaldocument[constructions-]{constructions}
\externaldocument[properties-]{properties}
\externaldocument[morphisms-]{morphisms}
\externaldocument[coherent-]{coherent}
\externaldocument[divisors-]{divisors}
\externaldocument[limits-]{limits}
\externaldocument[varieties-]{varieties}
\externaldocument[topologies-]{topologies}
\externaldocument[descent-]{descent}
\externaldocument[perfect-]{perfect}
\externaldocument[more-morphisms-]{more-morphisms}
\externaldocument[flat-]{flat}
\externaldocument[groupoids-]{groupoids}
\externaldocument[more-groupoids-]{more-groupoids}
\externaldocument[etale-]{etale}
\externaldocument[chow-]{chow}
\externaldocument[intersection-]{intersection}
\externaldocument[pic-]{pic}
\externaldocument[weil-]{weil}
\externaldocument[adequate-]{adequate}
\externaldocument[dualizing-]{dualizing}
\externaldocument[duality-]{duality}
\externaldocument[discriminant-]{discriminant}
\externaldocument[derham-]{derham}
\externaldocument[local-cohomology-]{local-cohomology}
\externaldocument[algebraization-]{algebraization}
\externaldocument[curves-]{curves}
\externaldocument[resolve-]{resolve}
\externaldocument[models-]{models}
\externaldocument[functors-]{functors}
\externaldocument[equiv-]{equiv}
\externaldocument[pione-]{pione}
\externaldocument[etale-cohomology-]{etale-cohomology}
\externaldocument[proetale-]{proetale}
\externaldocument[relative-cycles-]{relative-cycles}
\externaldocument[more-etale-]{more-etale}
\externaldocument[trace-]{trace}
\externaldocument[crystalline-]{crystalline}
\externaldocument[spaces-]{spaces}
\externaldocument[spaces-properties-]{spaces-properties}
\externaldocument[spaces-morphisms-]{spaces-morphisms}
\externaldocument[decent-spaces-]{decent-spaces}
\externaldocument[spaces-cohomology-]{spaces-cohomology}
\externaldocument[spaces-limits-]{spaces-limits}
\externaldocument[spaces-divisors-]{spaces-divisors}
\externaldocument[spaces-over-fields-]{spaces-over-fields}
\externaldocument[spaces-topologies-]{spaces-topologies}
\externaldocument[spaces-descent-]{spaces-descent}
\externaldocument[spaces-perfect-]{spaces-perfect}
\externaldocument[spaces-more-morphisms-]{spaces-more-morphisms}
\externaldocument[spaces-flat-]{spaces-flat}
\externaldocument[spaces-groupoids-]{spaces-groupoids}
\externaldocument[spaces-more-groupoids-]{spaces-more-groupoids}
\externaldocument[bootstrap-]{bootstrap}
\externaldocument[spaces-pushouts-]{spaces-pushouts}
\externaldocument[spaces-chow-]{spaces-chow}
\externaldocument[groupoids-quotients-]{groupoids-quotients}
\externaldocument[spaces-more-cohomology-]{spaces-more-cohomology}
\externaldocument[spaces-simplicial-]{spaces-simplicial}
\externaldocument[spaces-duality-]{spaces-duality}
\externaldocument[formal-spaces-]{formal-spaces}
\externaldocument[restricted-]{restricted}
\externaldocument[spaces-resolve-]{spaces-resolve}
\externaldocument[formal-defos-]{formal-defos}
\externaldocument[defos-]{defos}
\externaldocument[cotangent-]{cotangent}
\externaldocument[examples-defos-]{examples-defos}
\externaldocument[algebraic-]{algebraic}
\externaldocument[examples-stacks-]{examples-stacks}
\externaldocument[stacks-sheaves-]{stacks-sheaves}
\externaldocument[criteria-]{criteria}
\externaldocument[artin-]{artin}
\externaldocument[quot-]{quot}
\externaldocument[stacks-properties-]{stacks-properties}
\externaldocument[stacks-morphisms-]{stacks-morphisms}
\externaldocument[stacks-limits-]{stacks-limits}
\externaldocument[stacks-cohomology-]{stacks-cohomology}
\externaldocument[stacks-perfect-]{stacks-perfect}
\externaldocument[stacks-introduction-]{stacks-introduction}
\externaldocument[stacks-more-morphisms-]{stacks-more-morphisms}
\externaldocument[stacks-geometry-]{stacks-geometry}
\externaldocument[moduli-]{moduli}
\externaldocument[moduli-curves-]{moduli-curves}
\externaldocument[examples-]{examples}
\externaldocument[exercises-]{exercises}
\externaldocument[guide-]{guide}
\externaldocument[desirables-]{desirables}
\externaldocument[coding-]{coding}
\externaldocument[obsolete-]{obsolete}
\externaldocument[fdl-]{fdl}
\externaldocument[index-]{index}

% Theorem environments.
%
\theoremstyle{plain}
\newtheorem{theorem}[subsection]{Theorem}
\newtheorem{proposition}[subsection]{Proposition}
\newtheorem{lemma}[subsection]{Lemma}

\theoremstyle{definition}
\newtheorem{definition}[subsection]{Definition}
\newtheorem{example}[subsection]{Example}
\newtheorem{exercise}[subsection]{Exercise}
\newtheorem{situation}[subsection]{Situation}

\theoremstyle{remark}
\newtheorem{remark}[subsection]{Remark}
\newtheorem{remarks}[subsection]{Remarks}

\numberwithin{equation}{subsection}

% Macros
%
\def\lim{\mathop{\mathrm{lim}}\nolimits}
\def\colim{\mathop{\mathrm{colim}}\nolimits}
\def\Spec{\mathop{\mathrm{Spec}}}
\def\Hom{\mathop{\mathrm{Hom}}\nolimits}
\def\Ext{\mathop{\mathrm{Ext}}\nolimits}
\def\SheafHom{\mathop{\mathcal{H}\!\mathit{om}}\nolimits}
\def\SheafExt{\mathop{\mathcal{E}\!\mathit{xt}}\nolimits}
\def\Sch{\mathit{Sch}}
\def\Mor{\mathop{\mathrm{Mor}}\nolimits}
\def\Ob{\mathop{\mathrm{Ob}}\nolimits}
\def\Sh{\mathop{\mathit{Sh}}\nolimits}
\def\NL{\mathop{N\!L}\nolimits}
\def\CH{\mathop{\mathrm{CH}}\nolimits}
\def\proetale{{pro\text{-}\acute{e}tale}}
\def\etale{{\acute{e}tale}}
\def\QCoh{\mathit{QCoh}}
\def\Ker{\mathop{\mathrm{Ker}}}
\def\Im{\mathop{\mathrm{Im}}}
\def\Coker{\mathop{\mathrm{Coker}}}
\def\Coim{\mathop{\mathrm{Coim}}}

% Boxtimes
%
\DeclareMathSymbol{\boxtimes}{\mathbin}{AMSa}{"02}

%
% Macros for moduli stacks/spaces
%
\def\QCohstack{\mathcal{QC}\!\mathit{oh}}
\def\Cohstack{\mathcal{C}\!\mathit{oh}}
\def\Spacesstack{\mathcal{S}\!\mathit{paces}}
\def\Quotfunctor{\mathrm{Quot}}
\def\Hilbfunctor{\mathrm{Hilb}}
\def\Curvesstack{\mathcal{C}\!\mathit{urves}}
\def\Polarizedstack{\mathcal{P}\!\mathit{olarized}}
\def\Complexesstack{\mathcal{C}\!\mathit{omplexes}}
% \Pic is the operator that assigns to X its picard group, usage \Pic(X)
% \Picardstack_{X/B} denotes the Picard stack of X over B
% \Picardfunctor_{X/B} denotes the Picard functor of X over B
\def\Pic{\mathop{\mathrm{Pic}}\nolimits}
\def\Picardstack{\mathcal{P}\!\mathit{ic}}
\def\Picardfunctor{\mathrm{Pic}}
\def\Deformationcategory{\mathcal{D}\!\mathit{ef}}
\usepackage{cleveref}
\crefname{originalsection}{original section}{original sections}
\Crefname{originalsection}{Original section}{Original sections}
\makeatletter
\newlabel{cohomology-section-internal-hom}{{08DH}{source}{}{sourcepointer-1}{}}
\newlabel{cohomology-section-internal-hom@cref}{{[originalsection][0][]08DH}{[1][0][]source}{}{}{}}
\makeatother
\crefname{originallemma}{original lemma}{original lemmas}
\Crefname{originallemma}{Original lemma}{Original lemmas}
\makeatletter
\newlabel{lemma-affine-compare-bounded}{{06Z0}{source}{}{sourcepointer-2}{}}
\newlabel{lemma-affine-compare-bounded@cref}{{[originallemma][0][]06Z0}{[1][0][]source}{}{}{}}
\makeatother
\crefname{originallemma}{original lemma}{original lemmas}
\Crefname{originallemma}{Original lemma}{Original lemmas}
\makeatletter
\newlabel{lemma-pseudo-coherent-affine}{{08E7}{source}{}{sourcepointer-3}{}}
\newlabel{lemma-pseudo-coherent-affine@cref}{{[originallemma][0][]08E7}{[1][0][]source}{}{}{}}
\makeatother
\crefname{originallemma}{original lemma}{original lemmas}
\Crefname{originallemma}{Original lemma}{Original lemmas}
\makeatletter
\newlabel{lemma-perfect-affine}{{08EB}{source}{}{sourcepointer-4}{}}
\newlabel{lemma-perfect-affine@cref}{{[originallemma][0][]08EB}{[1][0][]source}{}{}{}}
\makeatother
\crefname{originallemma}{original lemma}{original lemmas}
\Crefname{originallemma}{Original lemma}{Original lemmas}
\makeatletter
\newlabel{cohomology-lemma-Rhom-complex-of-direct-summands-finite-free}{{08I5}{source}{}{sourcepointer-5}{}}
\newlabel{cohomology-lemma-Rhom-complex-of-direct-summands-finite-free@cref}{{[originallemma][0][]08I5}{[1][0][]source}{}{}{}}
\makeatother
\title{Quasi-coherence and affine computation of the derived internal Hom out of pseudo-coherent or perfect objects}
\author{The Stacks Project Authors}
\date{stacks-project:perfect.tex:local-snapshot-20261012; GNU FDL 1.2 or later}
\begin{document}
\maketitle
\setcounter{subsection}{1}
\noindent\textit{Editorial note.} All delivered mathematical text is the author's own, copied from the pinned Stacks Project chapter file perfect.tex; no statement or proof step was rewritten, reordered or removed. The counted span is the author environment and its own proof environment. References inside the span are handed to the assembler as navigation stubs that carry the canonical tag of the cited tag only; whether each of them is load-bearing is recorded on the stub itself.\par\medskip
\noindent\textit{Allowed standard prerequisites (editorial).}\begin{itemize}
\item The internal hom construction in the derived category of O-modules commutes with localization. invoked at proof lines 2687-2689. Reduces parts (1) and (2) to affine opens of X.
\item On affines the derived quasi-coherent categories compare with module categories, pseudo-coherent and perfect objects correspond to bounded above and finite complexes of finite projective modules. invoked at proof lines 2690-2693. Reduces (1) and (2) to the affine statement (3).
\item The internal hom of complexes represented by finite projective complexes is computed termwise, by the cohomology Rhom lemma for complexes of direct summands of finite free modules. invoked at proof lines 2695-2699. Proves part (3) in both the pseudo-coherent and the perfect case.
\item On a ringed space the nth cohomology sheaf of the internal hom is the sheaf associated to the presheaf of Ext groups on opens. invoked at proof lines 2702-2708. Gives the presheaf formula of part (4) on principal opens.
\end{itemize}\medskip
\ifdefined\hypertarget\hypertarget{sourcepointer-1}{}\fi
\noindent\textit{Source reference.} Original section 08DH: sources/stacks-chow/tags.txt, source lines 10637--10637. Cited at proof lines 2688-2689 and 2707-2708 for the localization and presheaf descriptions; the reductions are inside the counted span.\par\medskip
\ifdefined\hypertarget\hypertarget{sourcepointer-2}{}\fi
\noindent\textit{Source reference.} Original lemma 06Z0: sources/stacks-perfect/pinned-original-source.tex, source lines 454--454. Cited at proof lines 2690-2693 and 2714-2715 for the affine comparison; the reductions are inside the counted span.\par\medskip
\ifdefined\hypertarget\hypertarget{sourcepointer-3}{}\fi
\noindent\textit{Source reference.} Original lemma 08E7: sources/stacks-perfect/pinned-original-source.tex, source lines 2440--2440. Cited at proof lines 2691-2692 to reduce (1) to (3); the module computation is inside the counted span.\par\medskip
\ifdefined\hypertarget\hypertarget{sourcepointer-4}{}\fi
\noindent\textit{Source reference.} Original lemma 08EB: sources/stacks-perfect/pinned-original-source.tex, source lines 2636--2636. Cited at proof lines 2691-2693 to reduce (2) to (3); the module computation is inside the counted span.\par\medskip
\ifdefined\hypertarget\hypertarget{sourcepointer-5}{}\fi
\noindent\textit{Source reference.} Original lemma 08I5: sources/stacks-chow/tags.txt, source lines 10803--10803. Cited at proof lines 2696-2698 to prove part (3); the representation of K and L is inside the counted span.\par\medskip
\begin{lemma}
\label{lemma-quasi-coherence-internal-hom}
Let $X$ be a scheme.
\begin{enumerate}
\item If $L$ is in $D^+_\QCoh(\mathcal{O}_X)$ and
$K$ in $D(\mathcal{O}_X)$ is pseudo-coherent, then
$R\SheafHom(K, L)$ is in $D_\QCoh(\mathcal{O}_X)$
and locally bounded below.
\item If $L$ is in $D_\QCoh(\mathcal{O}_X)$ and
$K$ in $D(\mathcal{O}_X)$ is perfect, then
$R\SheafHom(K, L)$ is in $D_\QCoh(\mathcal{O}_X)$.
\item If $X = \Spec(A)$ is affine and $K, L \in D(A)$ then
$$
R\SheafHom(\widetilde{K}, \widetilde{L}) = \widetilde{R\Hom_A(K, L)}
$$
in the following two cases
\begin{enumerate}
\item $K$ is pseudo-coherent and $L$ is bounded below,
\item $K$ is perfect and $L$ arbitrary.
\end{enumerate}
\item If $X = \Spec(A)$ and $K, L$ are in $D(A)$, then the $n$th
cohomology sheaf of $R\SheafHom(\widetilde{K}, \widetilde{L})$
is the sheaf associated to the presheaf
$$
X \supset D(f) \longmapsto \Ext^n_{A_f}(K \otimes_A A_f, L \otimes_A A_f)
$$
for $f \in A$.
\end{enumerate}
\end{lemma}

\begin{proof}
The construction of the internal hom in the derived category of
$\mathcal{O}_X$ commutes with localization (see
Cohomology, Section \ref{cohomology-section-internal-hom}).
Hence to prove (1) and (2) we may replace $X$ by an affine open.
By Lemmas \ref{lemma-affine-compare-bounded},
\ref{lemma-pseudo-coherent-affine}, and
\ref{lemma-perfect-affine}
in order to prove (1) and (2) it suffices to prove (3).

\medskip\noindent
Part (3) follows from the computation of the
internal hom of Cohomology, Lemma
\ref{cohomology-lemma-Rhom-complex-of-direct-summands-finite-free}
by representing $K$ by a bounded above (resp.\ finite) complex of
finite projective $A$-modules and $L$ by a bounded below
(resp.\ arbitrary) complex of $A$-modules.

\medskip\noindent
To prove (4) recall that on any ringed space the $n$th cohomology sheaf of
$R\SheafHom(A, B)$ is the sheaf associated to the presheaf
$$
U \mapsto \Hom_{D(U)}(A|_U, B|_U[n]) =
\Ext^n_{D(\mathcal{O}_U)}(A|_U, B|_U)
$$
See Cohomology, Section \ref{cohomology-section-internal-hom}.
On the other hand, the restriction of $\widetilde{K}$ to a principal
open $D(f)$ is the image of $K \otimes_A A_f$ and similarly for $L$.
Hence (4) follows from the equivalence of categories of
Lemma \ref{lemma-affine-compare-bounded}.
\end{proof}

\end{document}
